% !TeX root = ../深度学习与大语言模型技术原理.tex
% 由主文档通过 \input 载入；请编译主文档。
\setcounter{chapter}{27} % 保留原书章号。
\chapter{GPU、CUDA、Triton 与编译优化}
\label{ch:kernels}

\section{并行层级与存储层级}

GPU 通过大量线程隐藏访存和计算延迟。CUDA 的线程（thread）组织为线程块（block），由流式多处理器（streaming multiprocessor, SM）调度；在 NVIDIA CUDA 架构中，线程束（warp）通常由32个线程组成。寄存器属于线程，片上共享内存可供同块线程协作，全局内存容量大但访问成本更高\cite{cuda}。

合并访存（memory coalescing）让相邻线程访问适合合并的地址，减少内存事务；非连续步长、错误布局和大量小访问会浪费带宽。共享内存分块允许一份数据被重复利用，但同步、银行冲突和寄存器压力也会影响性能。占用率（occupancy）表示可驻留活跃线程束比例，高占用率有时有助于隐藏延迟，却不是性能的唯一目标。

\section{Roofline 模型与算子融合}

设一个算子执行 $F$ 次浮点运算、搬运 $M$ 字节，算术强度（arithmetic intensity）为 $I=F/M$。Roofline 模型给出理想吞吐上界
\[
P_{\mathrm{achievable}}\le\min(P_{\mathrm{peak}},\mathcal B I),
\]
其中 $\mathcal B$ 为相关存储层带宽。若 $I$ 很小，增加算术峰值未必有用；若 $I$ 足够大，计算单元利用和数据类型成为重点。真实性能还受启动、同步和依赖影响，上式只是一种瓶颈分析框架。

算子融合（kernel fusion）把若干操作合成一个内核，避免中间张量反复写回显存。例如融合 RMSNorm 可在读取一行后完成平方归约、缩放与写出；融合 softmax 可避免分别调用最大值、指数、求和和除法内核。融合过多又可能增加寄存器占用、降低并行或限制形状适配，应实测端到端收益。

\section{Triton 内核与边界条件}

Triton 提供面向张量块的 GPU 编程方式，程序实例负责一块数据，通过显式索引、掩码、加载、计算与存储组织内核。它减少部分低层线程管理负担，但仍要求理解布局、块大小、归约和内存层级。官方 fused softmax 教程展示了通过融合减少访存的实例\cite{triton}。

以行 softmax 为例，一个程序读取若干行，在有效列上求最大值，计算指数和分母，再写回归一化结果。非整块长度需要掩码，填充位置应在最大值归约中用 $-\infty$，在指数求和后贡献零。全掩码行、极端数值、非连续输入和不同数据类型都应有明确支持范围。

验证自定义内核时，先与简单参考实现在随机、边界和极端输入上比较，再测性能。容差应结合数据类型与下游敏感性；一个内核在均匀随机输入上误差小，不保证训练中的离群激活也稳定。自动调优（autotuning）搜索块大小等配置，需要记录搜索开销与最终适用形状。

\section{计算图编译与性能测量}

LLM 编译器（LLM compiler）或通用张量编译器将模型图变换为更适合硬件的执行计划，包括常量折叠、算子融合、布局变换和内核选择。动态形状、Python 控制流和外部工具交互可能产生图中断或重新编译。CUDA Graph 一类机制减少重复启动开销，但对内存地址和执行结构有约束。

GPU 执行通常异步，直接在 CPU 端计时而不等待完成会低估耗时。应进行预热、同步或使用设备事件，区分编译与稳定运行成本，报告多次测量分布。优化局部内核后还需测端到端延迟，因为局部提速可能被通信、排队或其他算子掩盖。

